\documentclass[11pt,a4paper]{article}
\usepackage[utf8]{inputenc}
\usepackage[T1]{fontenc}
\usepackage[italian]{babel}
\usepackage{lmodern}
\usepackage[a4paper,margin=2.3cm]{geometry}
\usepackage{microtype}
\usepackage{amsmath}
\usepackage{booktabs}
\usepackage{tabularx}
\usepackage{xcolor}
\usepackage{listings}
\usepackage{hyperref}
\usepackage{enumitem}
\usepackage{fancyhdr}

\definecolor{accent}{HTML}{245B78}
\definecolor{codebg}{HTML}{F4F6F8}
\hypersetup{colorlinks=true,linkcolor=accent,urlcolor=accent}
\setlist{nosep,leftmargin=*}
\setlength{\parindent}{0pt}
\setlength{\parskip}{5pt}
\pagestyle{fancy}
\setlength{\headheight}{14pt}
\fancyhf{}
\lhead{ACO per CVRP}
\rhead{Analisi dei backend}
\cfoot{\thepage}
\lstset{basicstyle=\ttfamily\small,backgroundcolor=\color{codebg},frame=single,
  rulecolor=\color{codebg},breaklines=true,columns=fullflexible}

\title{\textbf{Implementazione ACO per il CVRP}\\
\large Backend sequenziale, OpenMP--MPI e CUDA}
\author{VehicleRoutingProblem --- documentazione tecnica}
\date{16 luglio 2026}

\begin{document}
\maketitle
\tableofcontents
\newpage

\section{Obiettivo e algoritmo comune}

Il progetto risolve il \emph{Capacitated Vehicle Routing Problem} (CVRP): dati un
deposito, $n$ clienti, $K$ veicoli e una capacità per veicolo, bisogna costruire
$K$ percorsi che partano e terminino al deposito, visitino tutti i clienti una
sola volta e minimizzino la distanza complessiva.

I tre backend impiegano \emph{Ant Colony Optimization} (ACO). Ogni formica
costruisce una soluzione scegliendo il prossimo cliente $j$ a partire dal nodo
$i$ con peso
\[
  w_{ij}=\tau_{ij}^{\alpha}\eta_{ij}^{\beta},
  \qquad \eta_{ij}=\frac{1}{d_{ij}},
\]
dove $\tau$ è il feromone, $\eta$ l'informazione euristica basata sulla distanza,
$\alpha$ controlla l'influenza del feromone e $\beta$ quella della distanza.
La scelta è probabilistica, tramite roulette-wheel selection.

Al termine di un'epoca si conserva la soluzione migliore, si applica
l'evaporazione
\[
  \tau_{ij}\leftarrow(1-\rho)\tau_{ij},
\]
e si deposita feromone sugli archi migliori, tipicamente in quantità
proporzionale a $Q/L$, dove $L$ è il costo della soluzione. I parametri usati
dalla CLI sono $\alpha=1$, $\beta=2$, $\rho=0.5$, $\tau_0=1$ e $Q=1$.

L'esecuzione può arrestarsi per timeout o per un numero configurabile di epoche
senza miglioramento significativo. Le direttive sono lette dalle variabili
\texttt{ACO\_SOLVER\_TIMEOUT\_SECONDS},
\texttt{ACO\_SOLVER\_STAGNATION\_EPOCHS} e
\texttt{ACO\_SOLVER\_MIN\_REL\_IMPROVEMENT}.

\section{Codice e librerie condivise}

Le implementazioni riusano i componenti in \texttt{src/common}:
\begin{itemize}
  \item \texttt{instance\_parser}: lettura delle istanze VRP/TSPLIB e delle coordinate;
  \item \texttt{matrix}: allocazione e rilascio della matrice delle distanze;
  \item \texttt{solution}: rappresentazione, copia e valutazione delle rotte;
  \item \texttt{aco\_shared}: generazione riproducibile dei seed e funzioni casuali.
\end{itemize}

Le librerie C principali sono \texttt{math.h} (collegata con \texttt{-lm}),
\texttt{stdint.h}, \texttt{float.h}, \texttt{stdlib.h} e \texttt{time.h}.
OpenMP aggiunge \texttt{omp.h}, MPI \texttt{mpi.h}; CUDA usa la runtime NVIDIA e
i tipi definiti dal compilatore \texttt{nvcc}.

\section{Backend sequenziale}

Il sorgente principale è \texttt{src/seq/aco\_sequential.c}. Tutte le formiche
sono costruite una dopo l'altra sul processore:
\begin{lstlisting}
per ogni epoca:
    aggiorna gli score tau^alpha * eta^beta
    per ogni formica:
        costruisci K rotte rispettando la capacita
        calcola il costo e conserva la migliore
    evapora e deposita il feromone
\end{lstlisting}

\subsection{Costruzione della soluzione}

Il workspace della formica contiene una soluzione riutilizzabile, il carico dei
veicoli, lo stato del generatore casuale e una bitmask dei clienti visitati. Un
cliente occupa un solo bit in un array di \texttt{uint64\_t}; il test avviene
con la parola \texttt{node/64} e il bit \texttt{node\%64}. Ciò riduce memoria e
traffico di cache rispetto a un array di interi.

La struttura chiamata candidate list include, in questa versione, tutti gli
$n$ clienti. Per ogni coppia utile vengono precalcolati l'indice del candidato e
$\eta^\beta$; a ogni epoca viene aggiornato lo score $\tau^\alpha\eta^\beta$.
La selezione scorre i clienti non visitati una prima volta per sommare gli score
e una seconda volta per applicare la roulette. Se non esiste un peso valido,
viene scelto il cliente non visitato più vicino.

Durante la costruzione viene preservata capacità sufficiente per i veicoli
successivi. In termini intuitivi, una rotta non assorbe clienti che saranno
necessari per rendere completabili le rotte rimanenti.

\subsection{Feromone e stagnazione}

Il deposito assegna il 30\% del peso alla migliore soluzione dell'epoca e il
70\% alla migliore globale. Il feromone è limitato tra \texttt{tau\_min} e
\texttt{tau\_max}, secondo un comportamento simile a Max--Min Ant System. Dopo
un periodo di stagnazione viene riavvicinato a $\tau_0$:
\[
  \tau_{ij}\leftarrow0.5\tau_{ij}+0.5\tau_0,
\]
così da aumentare nuovamente l'esplorazione.

\subsection{Ottimizzazioni CPU e AVX2}

Sono effettivamente presenti:
\begin{itemize}
  \item compilazione \texttt{-O3};
  \item memoria allineata a 64 byte e padding delle righe;
  \item puntatori \texttt{restrict}, utili all'analisi di aliasing del compilatore;
  \item $\eta^\beta$ e score in \texttt{float}, feromone e costo in \texttt{double};
  \item workspace riutilizzato, senza allocazioni nel ciclo delle formiche;
  \item casi rapidi per esponenti 1, 2 e 0.5, evitando \texttt{pow};
  \item bitset compatto per i clienti visitati.
\end{itemize}

Il codice contiene intrinsics AVX2 esplicite: \texttt{\_mm256\_i32gather\_pd}
carica quattro valori sparsi di feromone, \texttt{\_mm256\_mul\_pd} esegue le
moltiplicazioni e le conversioni trasformano i risultati tra \texttt{double} e
\texttt{float}. Il percorso SIMD viene compilato soltanto se è definita
\texttt{\_\_AVX2\_\_}. Il Makefile predefinito usa \texttt{-O3}, ma non abilita
automaticamente AVX2. Si può usare:
\begin{lstlisting}
make clean
make seq PERF_FLAGS="-mavx2 -mfma"
# oppure, per la CPU locale:
make seq PERF_FLAGS="-march=native"
\end{lstlisting}
\texttt{-march=native} può sfruttare ulteriori istruzioni della macchina, ma il
binario risultante è meno portabile. Inoltre un gather sparso è normalmente meno
efficiente di un caricamento vettoriale contiguo.

\section{Backend ibrido OpenMP--MPI}

Il sorgente è \texttt{src/openmp-mpi/aco\_parallel.c}. MPI distribuisce le
formiche tra processi, anche su nodi distinti; OpenMP utilizza i core condivisi
da ciascun rank. La compilazione usa \texttt{mpicc}, \texttt{-fopenmp},
\texttt{-DUSE\_MPI}, \texttt{-O3} e \texttt{-lm}.

MPI viene inizializzato con \texttt{MPI\_THREAD\_FUNNELED}: le chiamate MPI sono
eseguite dal thread master, mentre gli altri thread lavorano con OpenMP.

\subsection{Distribuzione MPI}

Le formiche sono divise quasi uniformemente tra i rank. Ogni rank conserva una
copia locale delle matrici, costruisce le proprie formiche con seed distinti e
calcola il proprio minimo. \texttt{MPI\_Allreduce} con \texttt{MPI\_MIN}
determina il miglior costo globale. Un altro Allreduce con \texttt{MPI\_MAX}
propaga la richiesta di arresto, impedendo che un rank continui mentre gli altri
sono terminati.

Il codice sincronizza i depositi come modifiche sparse, non scambiando tutta la
matrice $O(n^2)$. Ogni modifica contiene l'indice linearizzato dell'arco e un
incremento \texttt{float}. Un tipo MPI personalizzato viene costruito con
\texttt{MPI\_Type\_create\_struct}. I rank scambiano prima le cardinalità con
\texttt{MPI\_Allgather}, poi i dati con \texttt{MPI\_Iallgatherv}. La richiesta
non bloccante viene completata nell'epoca successiva con \texttt{MPI\_Wait}.
La quantità comunicata è proporzionale agli archi della soluzione, circa
$O(n+K)$, invece che agli elementi dell'intera matrice.

\subsection{Parallelismo OpenMP}

Una singola regione parallela persistente evita di ricreare il team a ogni
epoca. \texttt{proc\_bind(close)} tende a collocare vicini i thread, favorendo
la località di cache. Sono parallelizzati:
\begin{itemize}
  \item inizializzazione e conversione delle matrici;
  \item costruzione delle liste dei candidati;
  \item aggiornamento degli score;
  \item costruzione delle formiche;
  \item evaporazione della matrice;
  \item deposito sugli archi migliori.
\end{itemize}

I loop regolari usano \texttt{schedule(static)}. Le formiche usano
\texttt{schedule(runtime)}, configurabile tramite \texttt{OMP\_SCHEDULE}. Ogni
thread possiede un workspace privato. Una breve regione \texttt{critical}
riduce i migliori locali; \texttt{atomic} protegge aggiornamenti concorrenti del
feromone e l'inserimento nel buffer sparso. Le barriere separano le fasi che
dipendono dai risultati della fase precedente.

\subsection{Cache, candidate list e bitmask gerarchica}

Il programma legge, quando disponibile, la dimensione della cache L3 da Linux e
sceglie il numero di candidati per far occupare a indici ed euristiche circa il
70\% della L3:
\[
  k\approx\frac{0.7\,L3}{(n+1)\,8},\qquad16\leq k\leq512.
\]
Il working set passa così da $O(n^2)$ a $O(nk)$. Quando la lista locale non è
sufficiente viene eseguito un fallback globale.

I visitati sono organizzati su due livelli: L1 contiene un bit per cliente; L2
indica quali parole L1 hanno ancora clienti disponibili. Le istruzioni generate
da \texttt{\_\_builtin\_ctzll} individuano velocemente i bit attivi e permettono
di saltare blocchi completamente visitati.

Distanze, feromone ed euristiche sono prevalentemente \texttt{float}, con
matrici contigue, righe allineate a 64 byte e stride con padding. Questo dimezza
la memoria rispetto al \texttt{double}, migliora la cache e aumenta il possibile
throughput SIMD.

Nel backend principale non ci sono intrinsics AVX2 esplicite; i loop regolari
possono però essere autovettorizzati da GCC. Anche qui AVX2 non è garantito dal
Makefile di default e può essere richiesto con \texttt{PERF\_FLAGS=-mavx2} o
\texttt{-march=native}.

\section{Backend CUDA}

La logica host è in \path{src/cuda/aco_cuda.cu}, mentre i kernel sono in
\path{src/cuda/aco_cuda_kernels.cu}. La compilazione predefinita usa
\texttt{nvcc -O3 -std=c++17 -arch=sm\_75}; l'architettura può essere cambiata,
ad esempio con \texttt{make cuda CUDA\_ARCH=sm\_80}.

\subsection{Un warp per formica}

La scelta centrale è assegnare un warp di 32 thread a ogni formica. I lane del
warp collaborano per valutare 32 candidati, calcolare somme e prefissi, estrarre
il prossimo cliente ed eseguire il fallback globale.

Sono usate primitive warp-level:
\begin{itemize}
  \item \texttt{\_\_shfl\_down\_sync} per le riduzioni;
  \item \texttt{\_\_shfl\_up\_sync} per la somma prefissa;
  \item \texttt{\_\_shfl\_sync} per i broadcast;
  \item \texttt{\_\_ballot\_sync} e \texttt{\_\_ffs} per scegliere il lane valido;
  \item \texttt{\_\_syncwarp} per la sincronizzazione interna.
\end{itemize}
Queste operazioni scambiano dati attraverso i registri, senza richiedere shared
memory. La candidate list contiene 32 clienti, uno per lane, salvo istanze più
piccole.

\subsection{Feromone quantizzato e distanze on demand}

La matrice del feromone usa un solo \texttt{uint8\_t} per arco. Il byte codifica
logaritmicamente un valore tra $10^{-4}$ e $10^2$:
\[
  \tau=\exp(q\,\Delta+\log\tau_{\min}).
\]
Questo richiede un quarto della memoria di \texttt{float} e un ottavo di
\texttt{double}, aumentando il numero di valori trasferiti per transazione.
In scala logaritmica l'evaporazione moltiplicativa diventa approssimativamente
una sottrazione intera saturata, molto economica sull'intera matrice. Il costo è
una precisione inferiore e discretizzata.

La GPU non memorizza la matrice completa delle distanze: conserva le coordinate
come \texttt{float2} e calcola \texttt{sqrtf(dx*dx+dy*dy)} quando necessario.
Si scambia quindi calcolo, abbondante sulla GPU, con traffico e capacità di
memoria. La matrice quadratica del feromone rimane comunque allocata.

\subsection{Bitset, layout e atomiche}

Anche CUDA usa bitmask gerarchiche L1/L2 per i visitati. Le righe sono arrotondate
a 128 byte. I percorsi sono disposti come \texttt{[step][ant]}, così formiche
adiacenti accedono a indirizzi adiacenti; coordinate, candidate list e stati RNG
sono mantenuti in array dedicati. Queste scelte favoriscono accessi coalescenti.

Il deposito assegna un thread a ciascun veicolo. Poiché non esiste una normale
atomica per il singolo byte in questa forma, \texttt{atomicMax\_uint8} aggiorna
il byte con un ciclo \texttt{atomicCAS} sulla parola a 32 bit che lo contiene.

\subsection{Sequenza dei kernel e trasferimenti}

Ogni epoca esegue:
\begin{enumerate}
  \item reset dello stato delle formiche;
  \item costruzione parallela delle soluzioni;
  \item copia GPU--CPU dei sommari delle formiche;
  \item selezione CPU della migliore formica;
  \item eventuale copia del percorso migliore completo;
  \item evaporazione del feromone;
  \item deposito iteration-best con peso 30\%;
  \item deposito global-best con peso 70\%.
\end{enumerate}

Il ciclo host contiene varie \texttt{cudaDeviceSynchronize} e copie sincrone.
Questi punti possono incidere sulle istanze piccole. Riduzione del best sulla GPU,
stream asincroni, memoria host pinned e CUDA Graphs sarebbero possibili sviluppi,
ma non sono presenti nel codice corrente.

\section{Confronto conclusivo}

\small
\begin{tabularx}{\textwidth}{@{}lXXX@{}}
\toprule
\textbf{Aspetto} & \textbf{Sequenziale} & \textbf{OpenMP--MPI} & \textbf{CUDA}\\
\midrule
Parallelismo & Formiche seriali & Formiche tra rank e thread & Un warp per formica\\
Distanze & Matrice \texttt{double} & Matrice locale \texttt{float} & Coordinate \texttt{float2}, calcolo on demand\\
Feromone & Matrice \texttt{double} & Matrice \texttt{float} per rank & Matrice logaritmica \texttt{uint8\_t}\\
Candidati & Tutti gli $n$ & 16--512 adattati alla L3 & 32, uno per lane\\
Visitati & Bitset L1 & Bitset L1/L2 & Bitset L1/L2\\
SIMD & Intrinsics AVX2 condizionali & Autovettorizzazione possibile & Primitive SIMD del warp\\
Sincronizzazione & Nessuna tra worker & barrier, critical, atomic, collettive MPI & warp sync, atomiche, sync CPU--GPU\\
Comunicazione & Nessuna & Delta sparsi con \texttt{MPI\_Iallgatherv} & Copie CPU--GPU\\
\bottomrule
\end{tabularx}
\normalsize

In sintesi, il sequenziale è il riferimento CPU e include un percorso AVX2
esplicito ma non abilitato automaticamente. OpenMP--MPI scala distribuendo
formiche indipendenti e riduce il costo di cache e comunicazione con candidate
list adattive e delta sparsi. CUDA riorganizza invece dati e algoritmo intorno
all'hardware GPU: warp da 32 lane, shuffle nei registri, distanze calcolate al
momento e feromone compresso a 8 bit.

\section{Comandi di compilazione}
\begin{lstlisting}
make seq
make openmp_mpi
make cuda

# CPU locale con ISA disponibile sulla macchina
make clean
make seq PERF_FLAGS="-march=native"
make openmp_mpi PERF_FLAGS="-march=native"

# Target CUDA differente
make cuda CUDA_ARCH=sm_80
\end{lstlisting}

\end{document}
